\chapter{Protect the people}

Keys follow people. If the public identity, the family
WhatsApp, the daily-driver phone, and the signer are the
same blob, the architecture in the last chapter is
decoration.

The work here is shrinking what an attacker can learn,
and what they can grab that is standing next to you.

\section{Digital hygiene}

\textbf{Stop feeding OSINT.} Home address, plates, the
gym in the background, geotagged photos. Use throwaway
numbers and mail for things that do not need your real
identity. Full PII removal often needs a company or a
lawyer, not a browser extension. Decide what exposure you
will accept before you buy a cleanup service.

\textbf{Harden the devices that travel with you.}
Full-disk encryption. No seeds, no wallet apps on the
phone you take to dinner. No cloud sync of the sensitive
set. Factory-reset operational devices on a schedule if
they leave your sight.

\textbf{Do not glue KYC to the public handle.} Where the
law allows, keep the trading persona and the legal
identity from sitting in one screenshot.

\section{Split the identity}

Public posting and asset control should not share a
phone, an email, or a SIM.

A public phone for social and travel. A signer phone that
does not ride in the same pocket, and does not come to
the conference.

Family: no joint public profiles that map the household.
Separate devices. A simple rule that crypto holdings and
signing routines are not dinner conversation.

\textbf{Cold keys do not travel on your body} and they do
not live in a drawer on your route. Hardware and seed
backups sit in different places, none of which is ``the
apartment.''

\section{Family and staff}

Give the household one boring protocol they can actually
remember. Something like: if a stranger asks about money
or crypto, say nothing and send the pre-agreed word to
the pre-agreed person. Rehearse it once. Fancy duress
theatre that nobody remembers is not a protocol.

Drivers and assistants: vary routes, report being
followed, do not discuss assets or locations. Tabletop
the house and the office. NDAs do not replace drills.

Children get age-appropriate rules, not a lecture on
MPC. Door codes, pickup codes, an assembly point.

Panic hardware in a house or a car only helps if a
monitoring path exists and someone will answer it. Buying
the button is not the control.

\section{Routines}

Rotate the obvious loops. Do not make solo late-night
trips a habit if you are an obvious target. Ride-shares
do not need your legal name on the account the attacker
can already see. Do not carry a visible wad of cash
because you ``might need it.''

Home hardening in the source is basic: doors, windows,
cameras that actually alert someone, nothing that
advertises crypto on the shelf. This book will not pretend
that is a facility-security manual. It is the minimum so
that the rest of the design is not sitting behind a
cardboard door.

Travel: no signing on the trip if you can avoid it.
Someone trusted knows the itinerary. Devices in a safe
are still in a safe a hotel staffer can open. Skip the
high-profile side event if you are going alone and
angry-online.

\section{What you say}

Probes get a boring answer. ``I don't handle that.''
``Ask the company.'' Never confirm or deny holdings to a
curious stranger.

Cover stories have to be the same across family and
staff. ``It's in custody'' and ``the team runs treasury''
only work if nobody at the table contradicts them. If the
on-chain activity makes the story stupid, do not use it.

Once a quarter, look at what a stranger can already see:
profiles, address clustering, who on staff still has
access. Pay someone else to do it if you are the person
who posted the fundraising thread.
